% Official CC BY4.0 published HTML S2; original math alttext and author text, literal original numbering.


\par\medskip\noindent\textbf{2 Context and consequences of the main theorem}\par




\par\medskip\noindent\textbf{2.1 Complexity classes defined by invariant complexity measures}\par




A p-family is a sequence of polynomials $(f_{n})_{n\in\mathbb{N}}$ for which the number of variables of $f_{n}$ and the degree of $f_{n}$ are polynomially bounded functions of $n$. The class $\mathsf{VP}$ is the set of p-families whose algebraic circuit complexity is polynomially bounded.
Valiant’s landmark conjecture is that the permanent sequence $(\mathsf{per}_{n})_{n\in\mathbb{N}}$ with $\mathsf{per}_{n}=\sum_{\sigma\in\mathfrak{S}_{n}}\prod_{i=1}^{n}x_{i,\sigma(i)}$ is not contained in $\mathsf{VP}$, i.e., $\mathsf{VP}\neq\mathsf{VNP}$; for details see [14, 57, 13]. The extended conjecture states that $(\mathsf{per}_{n})\notin\mathsf{VQP}$ [12], i.e., that the algebraic circuit complexity of the permanent is not quasipolynomially bounded.





Let $\operatorname{cc}(f)$ denote the algebraic circuit complexity of the polynomial $f\in\mathbb{C}[x_{1},\ldots,x_{k}]$. It is easy to see that $\operatorname{cc}(f)=\operatorname{cc}(g\cdot f)$ for any $g\in\textnormal{GL}_{k}$: A circuit for $g\cdot f$ is obtained from the one for $f$ by applying $g$ at the input gates. Recall that in the definition of algebraic circuits we allow affine linear forms at the input gates. We say that the complexity measure $\operatorname{cc}$ is an invariant complexity measure.
In more restrictive circuit definitions, for instance if the input gates are required to be variables, the associated circuit complexity $c^{\prime}$ might change under the action of $\textnormal{GL}_{k}$. However, both definitions give the same class $\mathsf{VP}$.





Several other classical classes in algebraic complexity theory can be defined as the set of p-families for which some invariant complexity measure is polynomially bounded.
For example, the class $\mathsf{VF}$ is the set of p-families whose algebraic formula size is polynomially bounded; this is the smallest size of an algebraic circuit whose underlying directed graph is a tree. The class $\mathsf{VBP}$ is the set of p-families whose algebraic branching program width is polynomially bounded; we do not define this notion here, but it is an easy consequence of the definition that it is an invariant complexity measure, e.g. [68, Def. 2.2]. More classical classes can be readily defined in this way, for example the class of p-families of polynomially bounded circuit size and constant depth (still allowing arbitrary linear forms as inputs), such as $\Sigma\Pi\Sigma$, or the class of p-families of sums of polynomially many powering gates $\Sigma\Lambda\Sigma$.
Finally, we mention $\mathsf{VNP}$, which is the class of p-families of polynomially bounded permanental complexity: The smallest $r$ such that $f$ can be written as the permanent of an $r\times r$ matrix with affine linear entries. In all these cases, the relevant complexity measure is invariant.





Notable complexity measures that are not invariant are those related to the sparsity of the polynomial, in the sense of [49], or to restrictions on how often a variable can occur, such as read-once or read-$k$ models [27]. A non-invariant complexity measure $c^{\prime}$ can be converted into an invariant complexity measure $c$ by defining $c(f):=\min\{c^{\prime}(gf)\mid g\in\textnormal{GL}_{k}\}$, i.e., by taking the minimum complexity over all base changes, but for some non-invariant measures this can change the corresponding complexity class. In this paper, we only study complexity classes defined by invariant measures.






\par\medskip\noindent\textbf{2.2 Lower bounds via isotypic metapolynomials}\par




An important problem in algebraic complexity theory concerns proving concrete lower bounds for interesting complexity measures $c$. More precisely, given a polynomial $f_{\textup{hard}}$ and a certain integer $r$, one aims to show $c(f_{{\textup{hard}}})>r$ or equivalently $f_{{\textup{hard}}}\notin X_{r}$ where $X_{r}=\{f\in\mathbb{C}[x_{1},\ldots,x_{k}]_{d}\mid c(f)\leq r\}$ is the set of polynomials satisfying the upper bound $r$ on the complexity measure $c(\cdot)$. A standard approach is to determine a function $\Delta$ with the property that $\Delta(f)=0$ for every $f\in X_{r}$ and $\Delta(f_{\textup{hard}})\neq 0$.





Several lower bound methods in complexity theory are based on this approach, and most often the function $\Delta$ is a metapolynomial. This is the case for rank methods such as the method of partial derivatives [70, 22], the method of shifted partial derivatives [41, 28], and other augmented flattening methods such as Koszul–Young flattenings [54]: in all these cases the complexity lower bound is obtained in terms of a lower bound of the rank of a matrix whose entries depend on the metavariables, or equivalently on the nonvanishing of a suitable set of minors, which are metapolynomials. Rank methods yield the recent breakthrough of [55] giving the first superpolynomial lower bounds for explicit polynomials to be computed by low-product-depth algebraic circuits in large characteristic, further extended to any field in [30].
Similarly, the current best known lower bounds for the border determinantal complexity of the permanent, as well as the one of the sum of powers, are based on a non-degeneracy condition which can be translated into the vanishing of a suitable metapolynomial [53, 52]. Finally, the algebraic branching program lower bounds discussed in [50, 36] are built on geometric conditions of the zero set of the polynomial of interest, which in turn can be translated into the vanishing of metapolynomials.





When the complexity measure $c(\cdot)$ is invariant, one can use representation theory to enhance the search for the metapolynomials yielding a separation. Our main result Theorem 1.1 shows that this enhancement is not expensive: the computational overhead is quasi-polynomial.





To illustrate the effect of our result, first consider a simpler reduction. Let $\Delta$ be a metapolynomial vanishing on $X_{r}$ such that $\Delta(f_{\textup{hard}})\neq 0$. Write $\Delta=\sum_{i}\Delta^{(i)}$ for homogeneous metapolynomials $\Delta^{(i)}$ of degree $i$. If $X_{r}$ is invariant under rescaling, one has $\Delta^{(i)}(X_{r})=0$ for every $i$. On the other hand, there exists at least one $i$ such that $\Delta^{(i)}(f_{\textup{hard}})\neq 0$. The relevant homogeneous component $\Delta^{(i)}$ can be extracted from $\Delta$ with a simple interpolation argument, and one has $\operatorname{cc}(\Delta^{(i)})\leq(\deg(\Delta)+1)\cdot\operatorname{cc}(\Delta)$. In summary, if the lower bound $c(f_{\textup{hard}})>r$ can be proved via a metapolynomial $\Delta$ satisfying $\operatorname{cc}(\Delta)\leq s$, then the same lower bound can be proved via a homogeneous metapolynomial $\Delta^{(i)}$ satisfying $\operatorname{cc}(\Delta^{(i)})\leq(\deg(\Delta)+1)s$; and even further, the homogeneous metapolynomial $\Delta^{(i)}$ can be realized as the projection of $\Delta$ onto the $i$-th homogeneous component of the space of metapolynomials.





Theorem 1.1 shows similar results for other projections and in particular for the projection onto a specific isotypic component. Indeed, suppose $\Delta$ is a (homogeneous) metapolynomial vanishing on $X_{r}$ and such that $\Delta(f_{\textup{hard}})\neq 0$. Write $\Delta=\sum_{\lambda}\Delta^{(\lambda)}$ as the sum of its isotypic components. If $c(\cdot)$ is an invariant complexity measure, then $X_{r}$ is invariant under the action of $\textnormal{GL}_{k}$. In this case, one has $\Delta^{(\lambda)}(X_{r})=0$ for every $\lambda$. On the other hand, there exists at least one $\lambda$ such that $\Delta^{(\lambda)}(f_{\textup{hard}})\neq 0$. Theorem 1.1 provides an upper bound on the circuit complexity of $\Delta^{(\lambda)}$ in terms of the circuit complexity of $\Delta$, showing that if the lower bound $c(f_{\textup{hard}})>r$ can be proved via a metapolynomial $\Delta$ satisfying $\operatorname{cc}(\Delta)\leq s$, then the same lower bound can be proved via an isotypic metapolynomial with circuit complexity controlled by the parameter $s$. Theorem 1.1 further proves analogous results for more refined projections, such as the $T$-isotypic projections or the projection on the highest weight space. These are discussed in Section 2.3.






\par\medskip\noindent\textbf{2.3 Highest weight vectors and representation theoretic obstructions}\par




Theorem 1.1 can be used to convert lower bounds proved by $\Delta$ into lower bounds proved by a highest weight vector (HWV for short). To see this, note that if $\Delta$ vanishes on $X_{r}$ and $\Delta(f_{\textup{hard}})\neq 0$, then there exists a highest weight vector $\Delta^{\prime}$ vanishing on $X_{r}$ and such that $\Delta^{\prime}(g\cdot f_{\textup{hard}})\neq 0$ for some $g\in\textnormal{GL}_{k}$; in fact, this can be achieved with a random $g\in\textnormal{GL}_{k}$.
Hence a HWV proves the lower bound $c(g\cdot f_{{\textup{hard}}})>r$; if $c(\cdot)$ is an invariant complexity measure, the same lower bound holds for $f_{{\textup{hard}}}$. Similarly to before, Theorem 1.1 provides an upper bound on the circuit complexity $\operatorname{cc}(\Delta^{\prime})$ in terms of $\operatorname{cc}(\Delta)$, showing only a quasipolynomial blowup in the complexity of the metapolynomial.





A coarser view on HWVs is given via representation theoretic multiplicities, which is a central idea of geometric complexity theory [60, 61, 59].
Let $Y:=\overline{\textnormal{GL}_{k}\,f_{\textup{hard}}}$ denote the orbit closure, and let $X:=\overline{X_{s}}$.
It is easy to see that $f_{\textup{hard}}\in X$ if and only if $Y\subseteq X$.
The vector space of HWVs of weight $\lambda$ defines a vector space of polynomial functions on $X$; the dimension of such space is called
the multiplicity $\mathsf{mult}_{\lambda}\mathbb{C}[X]$ of $\lambda$ in the coordinate ring $\mathbb{C}[X]$. Standard facts in representation theory, and in particular Schur’s Lemma (see Lemma 4.4), guarantee that if $\mathsf{mult}_{\lambda}\mathbb{C}[X]>\mathsf{mult}_{\lambda}\mathbb{C}[Y]$ for some $\lambda$, then $X\not\subset Y$. Therefore, multiplicities give a method to potentially separate two varieties, and obtain a lower bound on $c(f_{\textup{hard}})$.
In this case, we say $\lambda$ is a multiplicity obstruction. If additionally
$\mathsf{mult}_{\lambda}\mathbb{C}[X]>0=\mathsf{mult}_{\lambda}\mathbb{C}[Y]$, then $\lambda$ is called an occurrence obstruction.
Occurrence obstructions [15, 16] and multiplicity obstructions [46] are sometimes sufficient to obtain separations. However, it is known that occurrence obstructions do not work in certain setups [47, 38, 18, 24].





Theorem 1.1 makes no statement about the viability of the method of multiplicity obstructions. Indeed, one main hope of geometric complexity theory is that information about $\mathsf{mult}_{\lambda}\mathbb{C}[X]$ can be obtained
without ever having to write down any circuit for a HWV.
Information can be gained
by decomposing much easier coordinate rings of orbits, see [19, 17]. This gives upper bounds on $\mathsf{mult}_{\lambda}\mathbb{C}[X]$ by using classical representation theoretic branching rules. The lower bounds on $\mathsf{mult}_{\lambda}\mathbb{C}[\overline{\textnormal{GL}_{k}f_{\textup{hard}}}]$ are more difficult to obtain, see [46] for a recent implementation.






\par\medskip\noindent\textbf{2.4 Border complexity}\par




Lower bounds achieved via metapolynomials are better described in the setting of border complexity. Generally, for a given algebraic complexity measure $c(\cdot)$, one can define a corresponding border measure, $\underline{c}(\cdot)$ as follows:



\[\begin{aligned}
\underline{c}(f)=\min\left\{r:\begin{array}[]{l}\text{there exists a sequence $f_{\varepsilon}$ such that $f=\lim_{\varepsilon\to 0}f_{\varepsilon}$}\\ \text{and $c(f_{\varepsilon})=r$ for all but finitely many $\varepsilon\in\mathbb{C}$}\end{array}\right\};
\end{aligned}\]


correspondingly, a complexity class $\mathcal{C}$ defined in terms of the growth of a given complexity measure naturally induces a border complexity class $\bar{\mathcal{C}}$, described by the growth of the corresponding border complexity measure. In fact, $\bar{\mathcal{C}}$ can be defined as the closure of $\mathcal{C}$ with respect to a suitable topology [48]. In the setting of algebraic circuits, the circuit size $\operatorname{cc}(\cdot)$ has a corresponding border circuit measure $\underline{\operatorname{cc}}(\cdot)$ which defines the complexity class $\overline{\mathsf{VP}}$.





A natural question concerns whether $\mathcal{C}=\bar{\mathcal{C}}$ and more generally how much larger $\bar{\mathcal{C}}$ is compared to $\mathcal{C}$. For example, the border class of sequences admitting (border) polynomial size $\Sigma\Lambda\Sigma$ circuits is contained in $\mathsf{VBP}$, see [29] and [9, Theorem 4.2]. In [25], the same containment was shown for the border classes of $\Sigma^{[k]}\Pi\Sigma$ circuits (for every constant $k$). In some restricted setting, for instance in the study of matrix multiplication complexity, it is known that a complexity measure and its associated border complexity are equivalent [7]. In general, the question is wide open for most complexity classes and, in particular, for $\mathsf{VP}$.





Metapolynomials characterize border complexity classes in the following sense: if a sequence of polynomials $(f_{n})$ does not belong to a border complexity class $\bar{\mathcal{C}}$ then there is a sequence of metapolynomials witnessing the non-membership. This is made precise in Theorem 2.2.






\par\medskip\noindent\textbf{2.5 Algebraic complexity and algebraic natural proofs}\par




Razborov and Rudich [66] introduced the notion of natural proofs in Boolean circuit complexity.





\par\medskip\noindent\textbf{Definition 2.1 (Natural property).}\par



A subset $P\subset\{f:\{0,1\}^{n}\rightarrow\{0,1\}\}$ of Boolean functions is said to be a natural property useful against a class $\mathcal{C}$ of Boolean circuits if the following is true:



\begin{itemize}


\item 
 


(Usefulness). Any Boolean function $f:\{0,1\}^{n}\rightarrow\{0,1\}$ that can be computed by a Boolean circuit in $\mathcal{C}$ does not have the property $P$.





\item 
 


(Constructivity). Given the truth table of a Boolean function $f:\{0,1\}^{n}\rightarrow\{0,1\}$, whether it has the property $P$ can be decided in time polynomial in the length of the input, i.e., in time $2^{O(n)}$.





\item 
 


(Largeness). For all large enough $n$, at least a $2^{-O(n)}$ fraction of all $n$ variate Boolean functions have the property $P$.





\end{itemize}





A proof that a certain family of Boolean functions cannot be computed by circuits in $\mathcal{C}$ is said to be a natural lower bound proof if the proof (perhaps implicitly) proceeds via establishing a natural property useful against $\mathcal{C}$, and showing that the candidate hard function has this property.
Razborov and Rudich showed that most of the known Boolean circuit lower bound proofs, such as lower bounds for $\mathsf{AC}^{0}$ (constant-depth) circuits have the natural property.
One can set the same framework in the algebraic setting; see [39, 40, 31, 21, 51]. The exact details of the definitions are not fully settled yet, see page 19 of [21].





Following [40], the stretched classes $\mathsf{metaVP}$ and $\mathsf{metaVQP}$ arise as follows.
The class $\mathsf{metaVP}$ is the class of sequences $(\Delta_{1},\Delta_{2},\ldots)$ of metapolynomials of format $(\delta(n),n,k(n))$ with $k(n)$ polynomially bounded in $n$, and $\delta(n)$ and $\operatorname{cc}(\Delta_{n})$ polynomially bounded in $N(n)=\binom{k(n)+n-1}{n}$, which is the number of degree $n$ monomials in $k(n)$ variables. The class $\mathsf{metaVQP}$ is defined analogously, but with $\operatorname{cc}(\Delta_{n})$ quasipolynomially bounded in $N(n)$.





We will now argue that $k(n)$ is bounded by a polynomial in $\log(N(n))$.
To see this, choose $\gamma$ such that $k(n)\leq n^{\gamma}$ for $n$ large enough.
Take any such large enough $n$ and make a case distinction.
If $k(n)\leq n$, then $N=\binom{k+n-1}{k-1}\geq(1+\frac{n}{k-1})^{k-1}$ implies $\log N\geq(k-1)\log(1+\frac{n}{k-1})\geq(k-1)$, since $\frac{n}{k-1}\geq 1$; here we use the fact that $\binom{m}{r}\geq(m/r)^{r}$.
On the other hand, if $k(n)\geq n+1$, then since $\binom{m}{r}$ is an increasing function on $m$ for a fixed $r$, we get that $N=\binom{k+n-1}{n}\geq\binom{2n}{n}\geq 2^{n}$, for $n\geq 1$.
Therefore, $(\log N)^{\gamma}\geq n^{\gamma}\geq k$.
Hence, the claim is proved in both cases.





Fix an invariant complexity measure $c$ and let $X_{d,r}$ denote the set of homogeneous degree $d$ polynomials $f$ with $c(f)\leq r$.
For a format $(\delta,d,k)$ metapolynomial $\Delta:\mathbb{C}[x_{1},\ldots,x_{k}]_{d}\to\mathbb{C}$ let $\Delta|_{X_{d,r}}:X_{d,r}\to\mathbb{C}$ denote the restriction of $\Delta$ to $X_{d,r}$.
Let $\mathcal{C}$ denote the set of p-families with polynomially bounded $c$.
The vanishing ideal of $\mathcal{C}$ is defined as





\[\begin{aligned}
\displaystyle I(\mathcal{C}):=\left\{(\Delta_{n})_{n\in\mathbb{N}}\textup{ of format $(\ast,n,\ast)$ }:\begin{array}[]{l}\text{for every polynomially bounded }r(n),\\ \text{the sequence }\big{(}\Delta_{1}|_{X_{1,r(1)}},\Delta_{2}|_{X_{2,r(2)}},\ldots\big{)}\\ \text{eventually vanishes identically}\end{array}\right\}
\end{aligned}\]




To emphasize how natural this definition is, we provide a corresponding Nullstellensatz.





\par\medskip\noindent\textbf{Theorem 2.2 (Nullstellensatz for $\mathcal{C}$).}\par



Let $(f_{n})_{n\in\mathbb{N}}$ be a p-family of homogeneous polynomials such that $(\deg f_{n})$ is a strictly increasing sequence. The sequence $(f_{n})$ lies in $\overline{\mathcal{C}}$ if and only if for every $(\Delta_{n})\in I(\mathcal{C})$ we have $\Delta_{\deg(f_{n})}(f_{n})=0$ eventually.






The proof of Theorem 2.2 is given in Appendix A.





\par\medskip\noindent\textbf{Definition 2.3 (Algebraic natural proof).}\par



We say that $\mathcal{C}$ has algebraic natural proofs if $I(\mathcal{C})\mathop{\cap}\mathsf{metaVQP}$ contains a sequence that is not eventually zero.






Here we loosened the definition from the literature slightly from $\mathsf{metaVP}$ (see [40, 21]) to $\mathsf{metaVQP}$, but note that the literature does not have just one standard definition, see page 19 of [21].
The landmark question in algebraic metacomplexity is whether $\mathsf{VP}$ has algebraic natural proofs.





Definition 2.3 should be compared with Definition 2.1 in the Boolean setting. If $I(\mathsf{VP})\mathop{\cap}\mathsf{metaVQP}$ contains an eventually nonzero sequence $\Delta$ of format $(\delta(n),n,k(n))$, then clearly there exists the largest $r(n)$ such that $\Delta_{n}(X_{n,r(n)})=\{0\}$, and further $\operatorname{cc}(\Delta_{n})\leq 2^{\operatorname{poly}(\log N)}$, where $N:=\binom{k(n)+n-1}{n}$, and $k(n)\leq\operatorname{poly}(n)$. Therefore, $\Delta$ certifies the non-membership, a useful natural property. Furthermore, it trivially satisfies the largeness criterion, since $\Delta_{n}$ does not vanish on most polynomials. Finally, the quasipolynomial bound on the algebraic circuit complexity of $\Delta_{n}$ gives a natural algebraic analogue of the notion of constructivity.





Let Isotypic denote the set of sequences of isotypic metapolynomials.
The following theorem follows from Theorem 1.1. It states in a concise way that only isotypic metapolynomials have to be considered for algebraic natural proofs.





\par\medskip\noindent\textbf{Theorem 2.4 (Main theorem about algebraic natural proofs).}\par



$\mathcal{C}$ has algebraic natural proofs if and only if $I(\mathcal{C})\mathop{\cap}\mathsf{metaVQP}\mathop{\cap}\textup{{Isotypic}}$ contains a sequence that is not eventually zero.






\par\medskip\noindent\textbf{Proof.}\par




One direction is clear, because $I(\mathcal{C})\mathop{\cap}\mathsf{metaVQP}\mathop{\cap}\textup{{Isotypic}}\subseteq I(\mathcal{C})\mathop{\cap}\mathsf{metaVQP}$.





For the other direction, we first consider a general property of isotypic components of vanishing ideals.
Let $I(X)_{\delta}=\{\Delta\textup{ of format $(\delta,d,k)$}\mid\Delta|_{X}=0\}$.
Let $X$ be closed under the action of $\textnormal{GL}_{k}$.
If $\Delta\in I(X)_{\delta}$, then $g\Delta\in I(X)_{\delta}$ for every $g\in\textnormal{GL}_{k}$.
Indeed, in this case for all $x\in X$ we have $(g\Delta)(x)=(\Delta)(g^{-1}x)=0$, since $g^{-1}x\in X$.
Moreover, $I(X)_{\delta}$ is closed under linear combinations, hence $I(X)_{\delta}$ is a vector space, and in particular closed under taking limits and thus closed under the action of the Lie algebra.
We have $\Delta\in I(X)_{\delta}$ if and only if for all $\lambda$ we have $\Delta^{(\lambda)}\in I(X)_{\delta}$, because the projection of $\Delta$ to an isotypic component $\Delta^{(\lambda)}$ involves only finite linear combinations of Lie algebra actions, see Section 5.





Let $\mathcal{C}$ have algebraic natural proofs, witnessed by a sequence $(\Delta_{n})_{n\in\mathbb{N}}\in I(\mathcal{C})\mathop{\cap}\mathsf{metaVQP}$ of format $(\delta(n),n,k(n))$ that is not eventually zero.
Now, let $(\lambda^{n})_{n\in\mathbb{N}}$ be a sequence of partitions, $\lambda^{n}\mathchoice{\scalebox{1.2}[0.7]{ $\vdash$}_{k(n)}\,}{\scalebox{1.2}[0.7]{ $\vdash$}_{k(n)}\,}{\scalebox{0.7}{\scalebox{1.2}[0.7]{$\vdash$}}_{k(n)}\,}{\scalebox{0.5}{\scalebox{1.2}[0.7]{$\vdash$}}_{k(n)}\,}{n\cdot\delta(n)}$, such that the isotypic component $\Delta_{n}^{(\lambda^{n})}$ of $\Delta_{n}$ is nonzero whenever $\Delta_{n}$ is nonzero.
By construction, $(\Delta_{n}^{(\lambda_{n})})_{n\in\mathbb{N}}\in\textup{{Isotypic}}$ and $(\Delta_{n}^{(\lambda_{n})})_{n\in\mathbb{N}}$ is not eventually zero.
Since each $X_{n,i}$ is invariant under the group action, we have for all $i$ and $n$:
Whenever $\Delta_{n}|_{X_{n,i}}=0$, then also $\Delta_{n}^{(\lambda^{n})}|_{X_{n,i}}=0$.
For every polynomially bounded $r$ we have
$\Delta_{n}|_{X_{n,r(n)}}=0$ eventually,
hence we have $\Delta_{n}^{(\lambda^{n})}|_{X_{n,r(n)}}=0$ eventually, and hence $(\Delta_{n}^{(\lambda_{n})})_{n\in\mathbb{N}}\in I(\mathcal{C})$.
Let $N(n)=\binom{k(n)+n-1}{n}$.
Let $s(n)=\operatorname{cc}(\Delta_{n})\in 2^{\operatorname{poly}(\log(N))}$.
By Theorem 1.1
with $k(n)\in\operatorname{poly}(\log(N))$, $\delta(n)\in\operatorname{poly}(N)$, and $d=n$, we
have $\operatorname{cc}(\Delta^{(\lambda^{n})}_{n})\in 2^{\operatorname{poly}(\log(N))}$,
which implies $(\Delta_{n}^{(\lambda_{n})})_{n\in\mathbb{N}}\in\mathsf{metaVQP}$,
which finishes the proof.
$\hfill\blacktriangleleft$





